\section{Background}
\label{sec:background}

\paragraph{Phases and steps.}
vLLM runs in \emph{engine steps}: each step, the scheduler picks the tokens to process
and the model runs one forward pass over all of them. A request first waits in a
\emph{queue}. Its \emph{prefill} phase processes the prompt and produces the first
output token. Its \emph{decode} phase then produces each remaining output token in one
step, so a request with 32 output tokens spends 31 steps in decode. With chunked
prefill \citep{agrawal2024sarathiserve}, a long prompt is split across several steps,
and one step can carry both decode tokens and prompt chunks.

\paragraph{Token budget.}
The \emph{token budget} $B$ (vLLM's \texttt{max\_num\_batched\_tokens}) is the most
tokens one step may process: one per decoding request plus every prompt token in the
step's chunks. Its default in our configuration is 2048. The scheduler fills each step
first with running requests (one token per decoding request, or the next chunk of a
prompt in progress) and then with waiting prompts, until the budget is used. A step
that carries a large chunk takes about as long as a prefill, and every request decoding
in that step waits for it. vLLM's documentation says smaller budgets give better
inter-token latency and larger ones better TTFT and throughput
\citep{vllmdocsoptimization}.

\paragraph{Metrics.}
Client metrics come from vLLM's benchmark. TTFT runs from sending a request to
receiving its first token, and E2E to receiving its last. Inter-token latency (ITL) is
each gap between consecutive tokens of a request; time per output token (TPOT) is a
request's average gap, $(\mathrm{E2E}-\mathrm{TTFT})/(n-1)$ for $n$ output tokens. We
report ITL percentiles, the median (p50) and 99th percentile (p99), because averages
hide jitter \citep{agrawal2024etalon}. The benchmark pools the gaps of all requests in a
run, and one step produces a gap for every request decoding in it, so a run's ITL p99
reflects its slowest few steps; where we give a mean ITL, it nearly equals the mean
TPOT. The \emph{share of latency before the first token} is $\sum \mathrm{TTFT} / \sum
\mathrm{E2E}$ over a run's requests. On the server, vLLM's histograms split each
request into queue, prefill (from first being scheduled to the first token), and
decode (from the first token to the last) time.
